\scanpage{30}
\deleted{qui décrivent $\unsure{\mathfrak S_3|\mathcal S_3}\times\mathbb Z/2$, on va relever ces générateurs en des éléments de $\mathcal E$, de façon que $l_\infty,l_1,l_0$ s'expriment comme des mots en ces éléments, et il suffit alors d'ajouter cette relation $l_\infty l_1l_0$ aux relations correspondant aux précédentes\ldots}

$\dot\sigma_0,\dot\sigma_1,\dot\sigma_\infty,\dot\rho,\tau$, avec les relations
\[
\left\{\begin{aligned}
\dot\sigma_0^2&=\dot\sigma_1^2=\dot\sigma_\infty^2=1\\
\dot\sigma_0\dot\sigma_1&=\dot\sigma_1\dot\sigma_\infty=\dot\sigma_\infty\dot\sigma_0=\dot\rho\\
\dot\rho^3&=1\\
\tau^2&=1\\
\tau&\text{ commute à }\dot\sigma_0,\dot\sigma_1,\dot\sigma_\infty\quad\text{(donc aussi à }\dot\rho\text{)}
\end{aligned}\right.
\]
on les relève en $\deleted{\tau_0,\tau_1,\tau_\infty}\ \sigma_0,\sigma_1,\sigma_\infty,\rho,\underbrace{\tau_0,\tau_1,\tau_\infty}_{\text{sur }\tau}$ avec les relations
\[
\left\{\begin{aligned}
\sigma_0^2&=\sigma_1^2=\sigma_\infty^2=\tau_0^2=[\tau_1^2=\tau_\infty^2]=1\\
\rho^3&=1\\
\rho\tau_0\rho^{-1}&=\tau_1,\quad\rho\tau_1\rho^{-1}=\tau_\infty,\quad\rho\tau_\infty\rho^{-1}=\tau_0\\
\tau_0\sigma_0&=\sigma_0\tau_0,\quad\tau_1\sigma_1=\sigma_1\tau_1,\quad\tau_\infty\sigma_\infty=\sigma_\infty\tau_\infty\\
\sigma_0\sigma_1&=\tau_\infty\tau_1\rho,\quad\sigma_1\sigma_\infty=\tau_0\tau_\infty\rho,\quad\sigma_\infty\sigma_0=\tau_1\tau_0\rho
\end{aligned}\right.
\]
\margin{NB Il faut ajouter les relations
\[
\left\{\begin{aligned}
\rho\sigma_0\rho^{-1}&=\sigma_1\\
\rho\sigma_1\rho^{-1}&=\sigma_\infty\\
\rho\sigma_\infty\rho^{-1}&=\sigma_0
\end{aligned}\right.
\]
\unsure{sont} conséquences des relations ci-contre}
\margin{Ce sont les éléments $\tilde\sigma_0,\tilde\sigma_1,\tilde\sigma_\infty$}
\margin{[NB On trouve
\[
\begin{gathered}
l_0=\tau_\infty\tau_1,\quad\deleted{\unsure{l_1=\tau_0\tau_\infty}},\\
l_1=\tau_0\deleted{\unsure{l_\infty=}}\tau_\infty\\
l_\infty=\tau_1\tau_0
\end{gathered}
\]
et $l_\infty l_1l_0=1$ est automatique]}
\margin{\deleted{\unsure{Ceci}} Ceci \unsure{implique} que $\tau_0,\tau_1,\tau_\infty$ \unsure{conjuguent|engendrent} \unsure{ss-groupe engendré} par $l_0,l_1,l_\infty$}

\deleted{\uline{Question} \deleted{\uline{Peut-on}}\inserted{On peut} exprimer $\rho$ en fonction de $\sigma_0,\sigma_1,\sigma_\infty,\tau_0,\tau_1,\tau_\infty$ ??, \deleted{\unsure{et pas plus}}\inserted{par l'une des} derniers groupes de relations, et substituer dans les autres.}
\mspage{29}
